\section{Metodologia}
La metodología empleada en esta práctica consistió en implementar el algoritmo de Gale--Shapley para resolver el Problema de Emparejamiento Estable y, posteriormente, evaluar experimentalmente su comportamiento temporal para distintos tamaños de entrada.

En una primera etapa se desarrolló una implementación eficiente del algoritmo considerando ambas variantes posibles: cuando el conjunto de hombres realiza las propuestas y cuando el conjunto de mujeres lo hace. Para optimizar el desempeño, se utilizó un esquema de \textit{ranking inverso} que permite comparar preferencias en tiempo constante $O(1)$, evitando recorridos lineales en las listas de preferencias.

En una segunda etapa se generaron automáticamente instancias de entrada para valores crecientes de $N$. Cada instancia contiene dos conjuntos de tamaño $N$ con listas completas de preferencias generadas aleatoriamente. La generación de datos se automatizó mediante un script independiente, lo que permitió controlar de manera sistemática el tamaño de las entradas y construir un conjunto de pruebas homogéneo.

Posteriormente, para cada instancia generada, se ejecutó el algoritmo y se registró el tiempo correspondiente al proceso de emparejamiento. Aunque en una implementación completa pueden distinguirse las fases de lectura, procesamiento e impresión, en este análisis se aisló el tiempo del algoritmo principal para que la medición reflejara de forma más precisa el costo computacional del procedimiento de Gale--Shapley.

Las mediciones se realizaron utilizando un temporizador de alta resolución (\texttt{perf\_counter}) y los resultados fueron almacenados en una tabla estructurada que relaciona el tamaño $N$ con el tiempo correspondiente.

Finalmente, los tiempos obtenidos fueron graficados utilizando tanto escala lineal como escala logarítmica para el eje de $N$, con el objetivo de analizar la forma del crecimiento temporal y contrastarlo con la complejidad teórica $O(N^2)$.

\subsection{Herramientas}

Para el desarrollo de la práctica se utilizaron las siguientes herramientas:

\begin{itemize}
    \item Lenguaje de programación Python 3 para la implementación del algoritmo y generación de instancias.
    \item Bibliotecas estándar de Python para medición de tiempos y manejo de estructuras de datos.
    \item Biblioteca \texttt{matplotlib} para la generación de gráficas.
    \item Entorno de desarrollo local y sistema de control de versiones Git para organización del código.
\end{itemize}
